\section{付加体の構成岩・形成過程}

\subsection{令和元年度再試験 問31}
\label{q:r01-ret-applied-31}

\noindent\textbf{原典：} sources/R01\_retest/17\_applied\_science.pdf、PDF物理10頁。

\textbf{公式正答：} 4（single）。

\paragraph{問題の要旨} 付加体へ取り込まれる典型的な岩石群から、最も関係の薄いものを選ぶ。

\Needspace{4\baselineskip}\paragraph{確認した条件・記述（独自要約）}\mbox{}\par

\noindent\textbf{方向：} 付加体の構成岩として最も関係のないもの。\par

\noindent\textbf{範囲：} 典型的な付加構成岩と後の貫入岩を区別。\par

\paragraph{問われる核心} 付加体の海洋プレート起源の岩石群

\paragraph{必要知識} 付加体の海洋プレート起源の岩石群。岩石名と生成過程、鉱物組合せ、温度・圧力・全岩組成、一般則と例外の区別。

\paragraph{考え方・正答を選ぶ理由（査読待ち初稿）} 海洋プレート本体・海山から玄武岩、遠洋の珪質堆積からチャート、海山上の炭酸塩堆積から石灰岩、海溝等の砕屑堆積から泥岩を対応させる。付加により取り込まれる典型的な岩石ではない花崗岩の4を選ぶ。現在の付加体地域に後で貫入した花崗岩が存在することとは両立する。

\Needspace{4\baselineskip}\paragraph{選択肢の検討（査読待ち初稿）}\mbox{}\par

\noindent\textbf{1：構成岩となり得る} 玄武岩。海洋地殻・海山などの火山岩が沈み込み時に付加され得る。\par

\noindent\textbf{2：構成岩となり得る} 泥岩。海溝などに運ばれる細粒砕屑堆積物が構成物となる。\par

\noindent\textbf{3：構成岩となり得る} チャート。海洋プレート上の遠洋性珪質堆積物として運搬・付加され得る。\par

\noindent\textbf{4：最も関係が薄い} 花崗岩。一般に地下の珪長質マグマの結晶化で形成され、海洋プレートの層序を剥ぎ取る付加物の典型ではない。付加後の花崗岩貫入まで否定しない。\par

\noindent\textbf{5：構成岩となり得る} 石灰岩。海山上の浅海生物起源の炭酸塩岩などが付加体へ取り込まれ得る。\par

\paragraph{再利用する考え方} 付加体の海洋プレート起源の岩石群を成因と観察される性質に結び付け、例示と全称命題、主要鉱物と副成分、公式解答と未解決の解釈を区別する。

\paragraph{共通理論} \hyperref[theory:petrology:accretion]{第\ref*{theory:petrology:accretion}節}

\par 参考：\cite{aist_accretion2026,gsj_saganoseki1993}。

\paragraph{関連問題} \hyperref[q:r05-reg-applied-26]{R05-REG-APPLIED-26}

\par 難易度 3（校準初稿）、優先度未評価。

\subsection{令和5年度 問26}
\label{q:r05-reg-applied-26}

\noindent\textbf{原典：} sources/R05/17\_applied\_science.pdf、PDF物理9頁。

\textbf{公式正答：} 2（single）。

\paragraph{問題の要旨} 付加体へ取り込まれる典型的な岩石群から、最も関係の薄いものを選ぶ。

\Needspace{4\baselineskip}\paragraph{確認した条件・記述（独自要約）}\mbox{}\par

\noindent\textbf{方向：} 付加体の構成岩として最も関係のないもの。\par

\noindent\textbf{範囲：} 典型的な付加構成岩と後の貫入岩を区別。\par

\paragraph{問われる核心} 付加体の海洋プレート起源の岩石群

\paragraph{必要知識} 付加体の海洋プレート起源の岩石群。岩石名と生成過程、鉱物組合せ、温度・圧力・全岩組成、一般則と例外の区別。

\paragraph{考え方・正答を選ぶ理由（査読待ち初稿）} 海洋プレート本体・海山から玄武岩、遠洋の珪質堆積からチャート、海山上の炭酸塩堆積から石灰岩、海溝等の砕屑堆積から泥岩を対応させる。付加により取り込まれる典型的な岩石ではない花崗岩の2を選ぶ。現在の付加体地域に後で貫入した花崗岩が存在することとは両立する。

\Needspace{4\baselineskip}\paragraph{選択肢の検討（査読待ち初稿）}\mbox{}\par

\noindent\textbf{1：構成岩となり得る} 玄武岩。海洋地殻・海山などの火山岩が沈み込み時に付加され得る。\par

\noindent\textbf{2：最も関係が薄い} 花崗岩。一般に地下の珪長質マグマの結晶化で形成され、海洋プレートの層序を剥ぎ取る付加物の典型ではない。付加後の花崗岩貫入まで否定しない。\par

\noindent\textbf{3：構成岩となり得る} 石灰岩。海山上の浅海生物起源の炭酸塩岩などが付加体へ取り込まれ得る。\par

\noindent\textbf{4：構成岩となり得る} チャート。海洋プレート上の遠洋性珪質堆積物として運搬・付加され得る。\par

\noindent\textbf{5：構成岩となり得る} 泥岩。海溝などに運ばれる細粒砕屑堆積物が構成物となる。\par

\paragraph{再利用する考え方} 付加体の海洋プレート起源の岩石群を成因と観察される性質に結び付け、例示と全称命題、主要鉱物と副成分、公式解答と未解決の解釈を区別する。

\paragraph{共通理論} \hyperref[theory:petrology:accretion]{第\ref*{theory:petrology:accretion}節}

\par 参考：\cite{aist_accretion2026,gsj_saganoseki1993}。

\paragraph{関連問題} \hyperref[q:r01-ret-applied-31]{R01-RET-APPLIED-31}

\par 難易度 3（校準初稿）、優先度未評価。

\section{火成岩組成・鉱物組合せ}

\subsection{平成27年度 問28}
\label{q:h27-reg-applied-28}

\noindent\textbf{原典：} sources/H27/17\_applied\_science.pdf、PDF物理11頁。

\textbf{公式正答：} 3（single）。

\paragraph{問題の要旨} 玄武岩の組成・起源・粘性について不適切な記述を選ぶ。

\Needspace{4\baselineskip}\paragraph{確認した条件・記述（独自要約）}\mbox{}\par

\noindent\textbf{方向：} 最も不適切なもの。\par

\noindent\textbf{対象：} 一般的な岩石分類・形成過程。\par

\paragraph{問われる核心} 玄武岩質マグマの組成・部分溶融・流動性

\paragraph{必要知識} 玄武岩質マグマの組成・部分溶融・流動性。岩石名と生成過程、鉱物組合せ、温度・圧力・全岩組成、一般則と例外の区別。

\paragraph{考え方・正答を選ぶ理由（査読待ち初稿）} 苦鉄質という分類がFe・Mgに富む組成と対応する。流紋岩はSiO2に富む珪長質岩であり、比較を逆にした3が誤りである。粘性と形状は組成だけでなく温度や結晶量にも依存するため、一般的傾向として読む。

\Needspace{4\baselineskip}\paragraph{選択肢の検討（査読待ち初稿）}\mbox{}\par

\noindent\textbf{1：適切（一般的分類）} 富士火山や伊豆大島を構成する苦鉄質火山岩である。。代表的な玄武岩火山の例である。ただし火山全体の全噴出物を単一岩種と断定しない。\par

\noindent\textbf{2：適切（一般的分類）} SiO2はおよそ50質量\%である。。玄武岩の代表的値で、全ての試料を正確に50\%とする意味ではない。\par

\noindent\textbf{3：不適切} 鉄・マグネシウムの含有量は流紋岩より少ない。。玄武岩は流紋岩より一般にFe・Mgに富むので、比較方向が逆。\par

\noindent\textbf{4：適切（一般的分類）} 上部マントル岩石の部分溶融が起源の一つである。。マントルのかんらん岩が部分溶融して玄武岩質メルトを生じ得る。起源を全溶融と混同しない。\par

\noindent\textbf{5：適切（一般的分類）} 流紋岩より粘性が低く、平らな火山体を作りやすい。。一般に高温でSiO2に乏しい玄武岩質溶岩は流紋岩質より流れやすい。噴火様式・供給量・地形も火山形状へ影響する。\par

\paragraph{再利用する考え方} 玄武岩質マグマの組成・部分溶融・流動性を成因と観察される性質に結び付け、例示と全称命題、主要鉱物と副成分、公式解答と未解決の解釈を区別する。

\paragraph{共通理論} \hyperref[theory:petrology:igneous]{第\ref*{theory:petrology:igneous}節}, \hyperref[theory:petrology:magma]{第\ref*{theory:petrology:magma}節}

\par 参考：\cite{jma_magma1999,gsj_igneous2021}。

\par 難易度 3（校準初稿）、優先度未評価。

\subsection{平成30年度 問30}
\label{q:h30-reg-applied-30}

\noindent\textbf{原典：} sources/H30/17\_applied\_science.pdf、PDF物理10頁。

\textbf{公式正答：} 5（single）。

\par\textbf{独立解答と公式正答に相違あり：} 独立解答 2|5。原典は花崗岩の主な有色鉱物に白雲母を含める。通常の無色分類と矛盾し、初回独立解答2|5、公式5。白雲母を改変しない。石英はんれい岩の例外と「主な」の範囲を含めAstraで査読。

\paragraph{問題の要旨} 火成岩の主要鉱物・斑晶と有色／無色の区別について不適切な記述を選ぶ。

\Needspace{4\baselineskip}\paragraph{確認した条件・記述（独自要約）}\mbox{}\par

\noindent\textbf{方向：} 最も不適切なもの。\par

\noindent\textbf{対象：} 一般的な岩石分類・形成過程。\par

\paragraph{問われる核心} 火成岩の鉱物組合せと有色分類の原典疑義

\paragraph{必要知識} 火成岩の鉱物組合せと有色分類の原典疑義。岩石名と生成過程、鉱物組合せ、温度・圧力・全岩組成、一般則と例外の区別。

\paragraph{考え方・正答を選ぶ理由（査読待ち初稿）} 一般的なはんれい岩へ石英を主要鉱物として入れた5が公式正答である。一方、原典に白雲母を「有色鉱物」とする記述もある。産総研の結晶構造ギャラリーは白雲母の色を無色と記載しており、初回独立解答は2|5であった。原典を黒雲母へ直したり白雲母も通常有色と教えたりせず、疑義を保存する。石英はんれい岩の例外も含め、公式の一般的主要鉱物の比較と厳密な分類を分ける。

\Needspace{4\baselineskip}\paragraph{選択肢の検討（査読待ち初稿）}\mbox{}\par

\noindent\textbf{1：適切（鉱物群の例示）} 花崗岩質岩石は石英・アルカリ長石・斜長石を無色鉱物として含む。。珪長質岩の主要な明色鉱物である。ただし鉱物量比が分類を左右し、あらゆる花崗岩類が同量の3種を含むとはしない。\par

\noindent\textbf{2：原典分類に疑義（通常は白雲母は無色）} 花崗岩質岩石の主な有色鉱物は白雲母・黒雲母・角閃石である。。黒雲母・角閃石は代表的な有色鉱物だが、白雲母は通常無色として扱う。白雲母が花崗岩に存在することと、有色分類が正しいことは別であり、原典のまま疑義を残す。\par

\noindent\textbf{3：適切（鉱物群の例示）} 玄武岩の斑晶には普通輝石・紫蘇輝石・斜長石・かんらん石がある。。玄武岩の種類により斑晶に現れ得る鉱物群である。全試料に全て同時に存在するという意味には読まない。紫蘇輝石は古い名称で斜方輝石の一種を指す。\par

\noindent\textbf{4：適切（鉱物群の例示）} 安山岩の斑晶には斜長石・普通輝石・紫蘇輝石・角閃石・磁鉄鉱・黒雲母・かんらん石がある。。安山岩の組成・水分・酸化状態等により現れ得る鉱物群である。かんらん石や黒雲母の存在の可否を、単一の岩種名だけで禁止しない。\par

\noindent\textbf{5：不適切（公式の選択対象）} はんれい岩の主な構成鉱物は普通輝石・斜長石・石英である。。一般的なはんれい岩は斜長石・輝石等を主成分とし、石英を典型的主要鉱物に並べる記述が公式の選択対象である。ただし石英はんれい岩という区分はあり、全はんれい岩で石英が絶対にないとはしない。\par

\paragraph{再利用する考え方} 火成岩の鉱物組合せと有色分類の原典疑義を成因と観察される性質に結び付け、例示と全称命題、主要鉱物と副成分、公式解答と未解決の解釈を区別する。

\paragraph{共通理論} \hyperref[theory:petrology:igneous]{第\ref*{theory:petrology:igneous}節}, \hyperref[theory:petrology:minerals]{第\ref*{theory:petrology:minerals}節}, \hyperref[theory:petrology:exceptions]{第\ref*{theory:petrology:exceptions}節}

\par 参考：\cite{aist_crystal_colors,gsj_igneous2021,usgs_rock_thesaurus}。

\paragraph{関連問題} \hyperref[q:r02-reg-applied-25]{R02-REG-APPLIED-25}

\par 難易度 3（校準初稿）、優先度未評価。

\subsection{令和2年度 問25}
\label{q:r02-reg-applied-25}

\noindent\textbf{原典：} sources/R02/17\_applied\_science.pdf、PDF物理9頁。

\textbf{公式正答：} 2（single）。

\par\textbf{独立解答と公式正答に相違あり：} 独立解答 1|2。原典は花崗岩の主な有色鉱物に白雲母を含める。通常の無色分類と矛盾し、初回独立解答1|2、公式2。白雲母を改変しない。石英はんれい岩の例外と「主な」の範囲を含めAstraで査読。

\paragraph{問題の要旨} 火成岩の主要鉱物・斑晶と有色／無色の区別について不適切な記述を選ぶ。

\Needspace{4\baselineskip}\paragraph{確認した条件・記述（独自要約）}\mbox{}\par

\noindent\textbf{方向：} 最も不適切なもの。\par

\noindent\textbf{対象：} 一般的な岩石分類・形成過程。\par

\paragraph{問われる核心} 火成岩の鉱物組合せと有色分類の原典疑義

\paragraph{必要知識} 火成岩の鉱物組合せと有色分類の原典疑義。岩石名と生成過程、鉱物組合せ、温度・圧力・全岩組成、一般則と例外の区別。

\paragraph{考え方・正答を選ぶ理由（査読待ち初稿）} 一般的なはんれい岩へ石英を主要鉱物として入れた2が公式正答である。一方、原典に白雲母を「有色鉱物」とする記述もある。産総研の結晶構造ギャラリーは白雲母の色を無色と記載しており、初回独立解答は1|2であった。原典を黒雲母へ直したり白雲母も通常有色と教えたりせず、疑義を保存する。石英はんれい岩の例外も含め、公式の一般的主要鉱物の比較と厳密な分類を分ける。

\Needspace{4\baselineskip}\paragraph{選択肢の検討（査読待ち初稿）}\mbox{}\par

\noindent\textbf{1：原典分類に疑義（通常は白雲母は無色）} 花崗岩質岩石の主な有色鉱物は白雲母・黒雲母・角閃石である。。黒雲母・角閃石は代表的な有色鉱物だが、白雲母は通常無色として扱う。白雲母が花崗岩に存在することと、有色分類が正しいことは別であり、原典のまま疑義を残す。\par

\noindent\textbf{2：不適切（公式の選択対象）} はんれい岩の主な構成鉱物は普通輝石・斜長石・石英である。。一般的なはんれい岩は斜長石・輝石等を主成分とし、石英を典型的主要鉱物に並べる記述が公式の選択対象である。ただし石英はんれい岩という区分はあり、全はんれい岩で石英が絶対にないとはしない。\par

\noindent\textbf{3：適切（鉱物群の例示）} 花崗岩質岩石は石英・アルカリ長石・斜長石を無色鉱物として含む。。珪長質岩の主要な明色鉱物である。ただし鉱物量比が分類を左右し、あらゆる花崗岩類が同量の3種を含むとはしない。\par

\noindent\textbf{4：適切（鉱物群の例示）} 玄武岩の斑晶には普通輝石・紫蘇輝石・斜長石・かんらん石がある。。玄武岩の種類により斑晶に現れ得る鉱物群である。全試料に全て同時に存在するという意味には読まない。紫蘇輝石は古い名称で斜方輝石の一種を指す。\par

\noindent\textbf{5：適切（鉱物群の例示）} 安山岩の斑晶には斜長石・普通輝石・紫蘇輝石・角閃石・磁鉄鉱・黒雲母・かんらん石がある。。安山岩の組成・水分・酸化状態等により現れ得る鉱物群である。かんらん石や黒雲母の存在の可否を、単一の岩種名だけで禁止しない。\par

\paragraph{再利用する考え方} 火成岩の鉱物組合せと有色分類の原典疑義を成因と観察される性質に結び付け、例示と全称命題、主要鉱物と副成分、公式解答と未解決の解釈を区別する。

\paragraph{共通理論} \hyperref[theory:petrology:igneous]{第\ref*{theory:petrology:igneous}節}, \hyperref[theory:petrology:minerals]{第\ref*{theory:petrology:minerals}節}, \hyperref[theory:petrology:exceptions]{第\ref*{theory:petrology:exceptions}節}

\par 参考：\cite{aist_crystal_colors,gsj_igneous2021,usgs_rock_thesaurus}。

\paragraph{関連問題} \hyperref[q:h30-reg-applied-30]{H30-REG-APPLIED-30}

\par 難易度 3（校準初稿）、優先度未評価。

\section{変成作用・鉱物組合せ・組織}

\subsection{令和元年度 問24}
\label{q:r01-reg-applied-24}

\noindent\textbf{原典：} sources/R01/17\_applied\_science.pdf、PDF物理8頁。

\textbf{公式正答：} 3（single）。

\paragraph{問題の要旨} 変成岩の形成要因、変成帯、高圧多形について不適切な記述を選ぶ。

\Needspace{4\baselineskip}\paragraph{確認した条件・記述（独自要約）}\mbox{}\par

\noindent\textbf{方向：} 最も不適切なもの。\par

\noindent\textbf{対象：} 一般的な岩石分類・形成過程。\par

\paragraph{問われる核心} 変成鉱物組合せの温度・圧力・組成依存と多形

\paragraph{必要知識} 変成鉱物組合せの温度・圧力・組成依存と多形。岩石名と生成過程、鉱物組合せ、温度・圧力・全岩組成、一般則と例外の区別。

\paragraph{考え方・正答を選ぶ理由（査読待ち初稿）} 鉱物組合せの平衡は温度・圧力の下で、供給される元素量を制約として自由エネルギーを最小化して決まる。したがって化学組成を無視する3が誤りである。石灰岩と泥岩を同じ温度圧力で変成しても同じ鉱物群になるとは限らない。

\Needspace{4\baselineskip}\paragraph{選択肢の検討（査読待ち初稿）}\mbox{}\par

\noindent\textbf{1：適切（一般的分類）} 火成岩・堆積岩・変成岩のいずれも原岩となる。。既存の変成岩が再び違う条件を受ける複変成もある。\par

\noindent\textbf{2：適切（一般的分類）} 固体中の新鉱物形成と外力による変形が変成作用の要素である。。再結晶・反応と変形を区別する。ただし接触変成などで顕著な変形を必須とはしない。\par

\noindent\textbf{3：不適切} 鉱物組合せは岩石の化学組成によらず温度・圧力のみで決まる。。温度・圧力に加えて全岩組成・流体組成等が必要である。元素が存在しなければその元素を要する鉱物はできない。\par

\noindent\textbf{4：適切（一般的分類）} 西南日本の高圧型変成帯は北から蓮華・周防・三波川の順に分布する。。原画像では蓮華・周防・三波川でありOCRの進華等を採用しない。高圧帯の典型的配置をいう。\par

\noindent\textbf{5：適切（一般的分類）} コース石・ダイヤモンドは石英・石墨のそれぞれ高圧多形である。。同じ化学組成でも結晶構造が違い、高圧側ではより稠密な構造が安定となる。コース石はSiO2、両炭素鉱物はCである。\par

\paragraph{再利用する考え方} 変成鉱物組合せの温度・圧力・組成依存と多形を成因と観察される性質に結び付け、例示と全称命題、主要鉱物と副成分、公式解答と未解決の解釈を区別する。

\paragraph{共通理論} \hyperref[theory:petrology:metamorphism]{第\ref*{theory:petrology:metamorphism}節}, \hyperref[theory:petrology:equilibrium]{第\ref*{theory:petrology:equilibrium}節}, \hyperref[theory:petrology:paired]{第\ref*{theory:petrology:paired}節}

\par 参考：\cite{nelson_assemblages2011,gsj_metamorphic_words,gsj_highpressure2016,hokkaido_polymorph2012}。

\paragraph{関連問題} \hyperref[q:r03-reg-applied-23]{R03-REG-APPLIED-23}

\par 難易度 3（校準初稿）、優先度未評価。

\subsection{令和3年度 問23}
\label{q:r03-reg-applied-23}

\noindent\textbf{原典：} sources/R03/17\_applied\_science.pdf、PDF物理8頁。

\textbf{公式正答：} 3（single）。

\paragraph{問題の要旨} 泥質変成岩の組織と変成度による岩型名を判定する。

\Needspace{4\baselineskip}\paragraph{確認した条件・記述（独自要約）}\mbox{}\par

\noindent\textbf{方向：} 最も不適切なもの。\par

\noindent\textbf{対象：} 一般的な岩石分類・形成過程。\par

\paragraph{問われる核心} 泥質岩の変成度と片理・片麻状組織

\paragraph{必要知識} 泥質岩の変成度と片理・片麻状組織。岩石名と生成過程、鉱物組合せ、温度・圧力・全岩組成、一般則と例外の区別。

\paragraph{考え方・正答を選ぶ理由（査読待ち初稿）} 粒の大きさと面・縞の構造を別々に確認する。片麻岩は粗粒という点だけで識別せず、鉱物の偏在による縞状組織が典型なので、それを否定する3が誤りである。

\Needspace{4\baselineskip}\paragraph{選択肢の検討（査読待ち初稿）}\mbox{}\par

\noindent\textbf{1：適切（一般的分類）} スレートは剥離性の発達した細粒岩で石材にも使われる。。細粒でもよく割れる劈開を持つ点が特徴で、屋根材などへ加工できる。\par

\noindent\textbf{2：適切（一般的分類）} 泥質岩の変成温度上昇に伴いスレート・千枚岩・片岩・片麻岩の順の変化が見られる。。泥質の代表的な進行変成系列である。岩型名は組成・変形・組織にも依存し、温度だけの万能温度計ではない。\par

\noindent\textbf{3：不適切} 片麻岩は粗粒で片状構造や縞状構造が見られない。。片麻岩は明暗の鉱物群が作る縞状組織などを持つ。見られないとの否定が誤り。\par

\noindent\textbf{4：適切（一般的分類）} 片岩は片理の発達した広域変成岩の総称である。。配向した板状鉱物等が面構造を作る。片理は単純な堆積時の層理とは異なる。\par

\noindent\textbf{5：適切（一般的分類）} ホルンフェルスは一般に片理・縞を持たず熱変成で形成される。。方向性に乏しい緻密な再結晶組織が典型である。原岩の構造が残ることまで絶対に排除しない。\par

\paragraph{再利用する考え方} 泥質岩の変成度と片理・片麻状組織を成因と観察される性質に結び付け、例示と全称命題、主要鉱物と副成分、公式解答と未解決の解釈を区別する。

\paragraph{共通理論} \hyperref[theory:petrology:textures]{第\ref*{theory:petrology:textures}節}, \hyperref[theory:petrology:metamorphism]{第\ref*{theory:petrology:metamorphism}節}

\par 参考：\cite{gsj_metamorphic_words,nelson_regional2011}。

\paragraph{関連問題} \hyperref[q:r01-reg-applied-24]{R01-REG-APPLIED-24}

\par 難易度 3（校準初稿）、優先度未評価。

\subsection{令和5年度 問24}
\label{q:r05-reg-applied-24}

\noindent\textbf{原典：} sources/R05/17\_applied\_science.pdf、PDF物理9頁。

\textbf{公式正答：} 5（single）。

\paragraph{問題の要旨} 固体の変成作用、接触・広域変成、弧–海溝系の温度圧力配置を判定する。

\Needspace{4\baselineskip}\paragraph{確認した条件・記述（独自要約）}\mbox{}\par

\noindent\textbf{方向：} 最も不適切なもの。\par

\noindent\textbf{対象：} 一般的な岩石分類・形成過程。\par

\paragraph{問われる核心} 変成過程と対の変成帯の熱・圧力の成因

\paragraph{必要知識} 変成過程と対の変成帯の熱・圧力の成因。岩石名と生成過程、鉱物組合せ、温度・圧力・全岩組成、一般則と例外の区別。

\paragraph{考え方・正答を選ぶ理由（査読待ち初稿）} 深く沈む比較的冷たいプレートは圧力の上昇に比べ温度上昇が小さい。火山弧下のマグマ近傍は熱源による高温が支配的になる。両側の低温高圧／高温低圧を逆にした5が誤りである。大理石は接触でも広域でもできるが、石灰岩の接触変成で形成されるという例示自体は正しい。

\Needspace{4\baselineskip}\paragraph{選択肢の検討（査読待ち初稿）}\mbox{}\par

\noindent\textbf{1：適切（一般的分類）} 岩石が高い温度・圧力の下で固体のまま鉱物や組織を変えると変成岩になる。。変成は固体状態の反応・再結晶・変形を中心とする。鉱物組成の変化と全岩化学組成の変化を区別し、全岩の外部物質交換を必須とはしない。\par

\noindent\textbf{2：適切（一般的分類）} マグマの貫入による周囲の加熱で接触変成岩を生じる。。局所的な熱源により貫入岩周辺に変成の帯を生じる。\par

\noindent\textbf{3：適切（一般的分類）} プレート運動に伴う広範囲の温度・圧力条件で広域変成岩を生じる。。地殻変動や沈み込みにより広い範囲で変成条件が成立する。高温と高圧が常に同じ程度である必要はない。\par

\noindent\textbf{4：適切（一般的分類）} 石灰岩が接触変成を受けると大理石となる。。炭酸塩鉱物の再結晶で大理石となる。広域変成によっても形成され、接触のみという排他的分類ではない。\par

\noindent\textbf{5：不適切} 弧–海溝系では海溝側に高温低圧型、陸側の火山前線下に低温高圧型ができる。。沈み込む冷たいプレート側は相対的な低温高圧、マグマによる熱源側は高温低圧なので逆である。\par

\paragraph{再利用する考え方} 変成過程と対の変成帯の熱・圧力の成因を成因と観察される性質に結び付け、例示と全称命題、主要鉱物と副成分、公式解答と未解決の解釈を区別する。

\paragraph{共通理論} \hyperref[theory:petrology:metamorphism]{第\ref*{theory:petrology:metamorphism}節}, \hyperref[theory:petrology:paired]{第\ref*{theory:petrology:paired}節}

\par 参考：\cite{aist_awa2021,nelson_contact2011,gsj_metamorphic_words}。

\paragraph{関連問題} \hyperref[q:r06-reg-applied-27]{R06-REG-APPLIED-27}, \hyperref[q:r07-reg-applied-23]{R07-REG-APPLIED-23}

\par 難易度 3（校準初稿）、優先度未評価。

\subsection{令和6年度 問27}
\label{q:r06-reg-applied-27}

\noindent\textbf{原典：} sources/R06/17\_applied\_science.pdf、PDF物理9頁。

\textbf{公式正答：} 4（single）。

\par\textbf{検討未確定：} 初回独立解答4と公式4は一致。5の「大理石は接触変成岩に属する」は代表例としては通用するが広域変成でも大理石が形成される。排他的な全称分類か例示かの読みをAstraで確認。

\paragraph{問題の要旨} 変成作用の定義と領家・三波川の変成型を判定する。

\Needspace{4\baselineskip}\paragraph{確認した条件・記述（独自要約）}\mbox{}\par

\noindent\textbf{方向：} 最も不適切なもの。\par

\noindent\textbf{対象：} 一般的な岩石分類・形成過程。\par

\paragraph{問われる核心} 領家・三波川の変成型と大理石分類の適用範囲

\paragraph{必要知識} 領家・三波川の変成型と大理石分類の適用範囲。岩石名と生成過程、鉱物組合せ、温度・圧力・全岩組成、一般則と例外の区別。

\paragraph{考え方・正答を選ぶ理由（査読待ち初稿）} 領家はマグマに関連する高温低圧型、三波川は沈み込みに関連する低温高圧型である。逆にした4が公式正答かつ初回独立解答である。ただし5は大理石の広域変成での形成を省いており、全称的な分類として読む場合には追加の疑義がある。原典の断定をそのまま普遍則として教えず、解釈確認を残す。

\Needspace{4\baselineskip}\paragraph{選択肢の検討（査読待ち初稿）}\mbox{}\par

\noindent\textbf{1：適切} 岩石が高い温度・圧力の下で固体のまま鉱物や組織を変えると変成岩になる。。変成は固体状態の反応・再結晶・変形を中心とする。鉱物組成の変化と全岩化学組成の変化を区別し、全岩の外部物質交換を必須とはしない。\par

\noindent\textbf{2：適切} プレート運動に伴う広範囲の温度・圧力条件で広域変成岩を生じる。。地殻変動や沈み込みにより広い範囲で変成条件が成立する。高温と高圧が常に同じ程度である必要はない。\par

\noindent\textbf{3：適切} マグマの貫入による周囲の加熱で接触変成岩を生じる。。局所的な熱源により貫入岩周辺に変成の帯を生じる。\par

\noindent\textbf{4：不適切} 領家帯は低温高圧、三波川帯は高温低圧である。。領家は高温低圧、三波川は低温高圧の代表であり、対応が逆。\par

\noindent\textbf{5：代表例として適切／全称解釈に疑義} 結晶片岩・片麻岩は広域変成岩、ホルンフェルス・大理石は接触変成岩に属する。。試験では代表的な形成例の対応としているが、大理石は石灰質原岩の広域変成でもできる。全ての大理石を接触起源と断定するなら誤りとなるので、例示と排他的区分を区別する。\par

\paragraph{再利用する考え方} 領家・三波川の変成型と大理石分類の適用範囲を成因と観察される性質に結び付け、例示と全称命題、主要鉱物と副成分、公式解答と未解決の解釈を区別する。

\paragraph{共通理論} \hyperref[theory:petrology:paired]{第\ref*{theory:petrology:paired}節}, \hyperref[theory:petrology:exceptions]{第\ref*{theory:petrology:exceptions}節}

\par 参考：\cite{aist_awa2021,nelson_regional2011,nelson_contact2011}。

\paragraph{関連問題} \hyperref[q:r05-reg-applied-24]{R05-REG-APPLIED-24}

\par 難易度 3（校準初稿）、優先度未評価。

\subsection{令和7年度 問23}
\label{q:r07-reg-applied-23}

\noindent\textbf{原典：} sources/R07/17\_applied\_science.pdf、PDF物理8頁。

\textbf{公式正答：} 4（single）。

\paragraph{問題の要旨} 固体の変成作用、接触・広域変成、弧–海溝系の温度圧力配置を判定する。

\Needspace{4\baselineskip}\paragraph{確認した条件・記述（独自要約）}\mbox{}\par

\noindent\textbf{方向：} 最も不適切なもの。\par

\noindent\textbf{対象：} 一般的な岩石分類・形成過程。\par

\paragraph{問われる核心} 変成過程と対の変成帯の熱・圧力の成因

\paragraph{必要知識} 変成過程と対の変成帯の熱・圧力の成因。岩石名と生成過程、鉱物組合せ、温度・圧力・全岩組成、一般則と例外の区別。

\paragraph{考え方・正答を選ぶ理由（査読待ち初稿）} 深く沈む比較的冷たいプレートは圧力の上昇に比べ温度上昇が小さい。火山弧下のマグマ近傍は熱源による高温が支配的になる。両側の低温高圧／高温低圧を逆にした4が誤りである。大理石は接触でも広域でもできるが、石灰岩の接触変成で形成されるという例示自体は正しい。

\Needspace{4\baselineskip}\paragraph{選択肢の検討（査読待ち初稿）}\mbox{}\par

\noindent\textbf{1：適切（一般的分類）} 岩石が高い温度・圧力の下で固体のまま鉱物や組織を変えると変成岩になる。。変成は固体状態の反応・再結晶・変形を中心とする。鉱物組成の変化と全岩化学組成の変化を区別し、全岩の外部物質交換を必須とはしない。\par

\noindent\textbf{2：適切（一般的分類）} マグマの貫入による周囲の加熱で接触変成岩を生じる。。局所的な熱源により貫入岩周辺に変成の帯を生じる。\par

\noindent\textbf{3：適切（一般的分類）} プレート運動に伴う広範囲の温度・圧力条件で広域変成岩を生じる。。地殻変動や沈み込みにより広い範囲で変成条件が成立する。高温と高圧が常に同じ程度である必要はない。\par

\noindent\textbf{4：不適切} 弧–海溝系では海溝側に高温低圧型、陸側の火山前線下に低温高圧型ができる。。沈み込む冷たいプレート側は相対的な低温高圧、マグマによる熱源側は高温低圧なので逆である。\par

\noindent\textbf{5：適切（一般的分類）} 石灰岩が接触変成を受けると大理石となる。。炭酸塩鉱物の再結晶で大理石となる。広域変成によっても形成され、接触のみという排他的分類ではない。\par

\paragraph{再利用する考え方} 変成過程と対の変成帯の熱・圧力の成因を成因と観察される性質に結び付け、例示と全称命題、主要鉱物と副成分、公式解答と未解決の解釈を区別する。

\paragraph{共通理論} \hyperref[theory:petrology:metamorphism]{第\ref*{theory:petrology:metamorphism}節}, \hyperref[theory:petrology:paired]{第\ref*{theory:petrology:paired}節}

\par 参考：\cite{aist_awa2021,nelson_contact2011,gsj_metamorphic_words}。

\paragraph{関連問題} \hyperref[q:r05-reg-applied-24]{R05-REG-APPLIED-24}

\par 難易度 3（校準初稿）、優先度未評価。

\section{火山砕屑物・堆積構造}

\subsection{令和3年度 問24}
\label{q:r03-reg-applied-24}

\noindent\textbf{原典：} sources/R03/17\_applied\_science.pdf、PDF物理9頁。

\textbf{公式正答：} 4（single）。

\paragraph{問題の要旨} 火山砕屑物の名称・粒径・堆積様式・溶結について不適切な記述を選ぶ。

\Needspace{4\baselineskip}\paragraph{確認した条件・記述（独自要約）}\mbox{}\par

\noindent\textbf{方向：} 最も不適切なもの。\par

\noindent\textbf{対象：} 一般的な岩石分類・形成過程。\par

\paragraph{問われる核心} 火山砕屑物の粒径・支持構造・溶結

\paragraph{必要知識} 火山砕屑物の粒径・支持構造・溶結。岩石名と生成過程、鉱物組合せ、温度・圧力・全岩組成、一般則と例外の区別。

\paragraph{考え方・正答を選ぶ理由（査読待ち初稿）} 名称、粒径、支持構造、堆積後の結合を別軸で比較する。降下軽石の比較的よくそろった粒子支持と、火砕流の灰を含む基質支持を入れ替えた4が不適切である。火山弾を火山岩塊より必ず大粒とする粒径順序へは読み替えない。

\Needspace{4\baselineskip}\paragraph{選択肢の検討（査読待ち初稿）}\mbox{}\par

\noindent\textbf{1：適切（一般的分類）} 多孔質で低い見かけ密度の白〜灰色砕屑物は軽石、暗色のものはスコリアという。。発泡したマグマ由来の砕屑物で、色・組成に差がある。低い見かけ密度と鉱物そのものの密度を混同しない。\par

\noindent\textbf{2：適切（一般的分類）} 粒径区分は細かい方から火山灰・火山礫・火山岩塊および火山弾である。。灰は2 mm未満、火山礫は2〜64 mm、岩塊・火山弾は64 mm超が基準である。岩塊と火山弾は同じ粗粒区分で、主に固結状態・形状で分ける。\par

\noindent\textbf{3：適切（一般的分類）} アグロメレートは火山灰基質に火山弾・本質岩塊を多く含み火口近くに分布する。。火口付近に堆積しやすい大きな本質砕屑物を多く含む火砕岩の説明である。\par

\noindent\textbf{4：不適切} 降下軽石は基質支持、火砕流堆積物は岩片支持である。。典型的な降下軽石は比較的淘汰され粒子同士で支持し、火砕流は細粒灰を含む不淘汰な基質支持が一般的であり、比較方向が逆。流れの相や堆積位置による例外はある。\par

\noindent\textbf{5：適切（一般的分類）} アグルチネートは放出された高温の岩片が溶結したものである。。まだ高温で粘着性のある放出物が集積し結合した岩石である。冷たい砕屑物の機械的固結だけとは異なる。\par

\paragraph{再利用する考え方} 火山砕屑物の粒径・支持構造・溶結を成因と観察される性質に結び付け、例示と全称命題、主要鉱物と副成分、公式解答と未解決の解釈を区別する。

\paragraph{共通理論} \hyperref[theory:petrology:pyroclast]{第\ref*{theory:petrology:pyroclast}節}, \hyperref[theory:petrology:deposition]{第\ref*{theory:petrology:deposition}節}

\par 参考：\cite{usgs_tephra_sizes,usgs_pinatubo1996,usgs_ash_density,usgs_bishop_tuff}。

\paragraph{関連問題} \hyperref[q:r06-reg-applied-26]{R06-REG-APPLIED-26}

\par 難易度 3（校準初稿）、優先度未評価。

\subsection{令和6年度 問26}
\label{q:r06-reg-applied-26}

\noindent\textbf{原典：} sources/R06/17\_applied\_science.pdf、PDF物理9頁。

\textbf{公式正答：} 5（single）。

\paragraph{問題の要旨} 火山砕屑物の名称・粒径・堆積様式・溶結について不適切な記述を選ぶ。

\Needspace{4\baselineskip}\paragraph{確認した条件・記述（独自要約）}\mbox{}\par

\noindent\textbf{方向：} 最も不適切なもの。\par

\noindent\textbf{対象：} 一般的な岩石分類・形成過程。\par

\paragraph{問われる核心} 火山砕屑物の粒径・支持構造・溶結

\paragraph{必要知識} 火山砕屑物の粒径・支持構造・溶結。岩石名と生成過程、鉱物組合せ、温度・圧力・全岩組成、一般則と例外の区別。

\paragraph{考え方・正答を選ぶ理由（査読待ち初稿）} 名称、粒径、支持構造、堆積後の結合を別軸で比較する。降下軽石の比較的よくそろった粒子支持と、火砕流の灰を含む基質支持を入れ替えた5が不適切である。火山弾を火山岩塊より必ず大粒とする粒径順序へは読み替えない。

\Needspace{4\baselineskip}\paragraph{選択肢の検討（査読待ち初稿）}\mbox{}\par

\noindent\textbf{1：適切（一般的分類）} 多孔質で低い見かけ密度の白〜灰色砕屑物は軽石、暗色のものはスコリアという。。発泡したマグマ由来の砕屑物で、色・組成に差がある。低い見かけ密度と鉱物そのものの密度を混同しない。\par

\noindent\textbf{2：適切（一般的分類）} 粒径区分は細かい方から火山灰・火山礫・火山岩塊および火山弾である。。灰は2 mm未満、火山礫は2〜64 mm、岩塊・火山弾は64 mm超が基準である。岩塊と火山弾は同じ粗粒区分で、主に固結状態・形状で分ける。\par

\noindent\textbf{3：適切（一般的分類）} アグロメレートは火山灰基質に火山弾・本質岩塊を多く含み火口近くに分布する。。火口付近に堆積しやすい大きな本質砕屑物を多く含む火砕岩の説明である。\par

\noindent\textbf{4：適切（一般的分類）} アグルチネートは放出された高温の岩片が溶結したものである。。まだ高温で粘着性のある放出物が集積し結合した岩石である。冷たい砕屑物の機械的固結だけとは異なる。\par

\noindent\textbf{5：不適切} 降下軽石は基質支持、火砕流堆積物は岩片支持である。。典型的な降下軽石は比較的淘汰され粒子同士で支持し、火砕流は細粒灰を含む不淘汰な基質支持が一般的であり、比較方向が逆。流れの相や堆積位置による例外はある。\par

\paragraph{再利用する考え方} 火山砕屑物の粒径・支持構造・溶結を成因と観察される性質に結び付け、例示と全称命題、主要鉱物と副成分、公式解答と未解決の解釈を区別する。

\paragraph{共通理論} \hyperref[theory:petrology:pyroclast]{第\ref*{theory:petrology:pyroclast}節}, \hyperref[theory:petrology:deposition]{第\ref*{theory:petrology:deposition}節}

\par 参考：\cite{usgs_tephra_sizes,usgs_pinatubo1996,usgs_ash_density,usgs_bishop_tuff}。

\paragraph{関連問題} \hyperref[q:r03-reg-applied-24]{R03-REG-APPLIED-24}

\par 難易度 3（校準初稿）、優先度未評価。
